
% Estructura del proyecto:
%
% main.tex :  El documento principal donde está la configuración y donde se insertan el resto de documentos.
% main.bib :  Archivo donde está la bibliografía en formato .bib
% chapters/. : La carpeta donde están los archivos .tex de los capítulos que insertamos en main.tex con el comando \input{chapters/capitulo}.
% documents/. : La carpeta donde están algunos documentos como la Cubierta, la Portada, la Acreditación del supervisor y la Contraportada.
% images/. : La carpeta donde están todas las imágenes.

%------------------- Configuración -------------------------------------
\documentclass[12pt,table]{book}
\usepackage[utf8]{inputenc}
\usepackage[english]{babel} % Idioma, inglés por defecto
% Ejemplo con formato B5 (no es obligatorio), ajustar márgenes y formato a conveniencia
\usepackage[b5paper,top=40mm, inner=30mm, outer=20mm, bottom=30mm]{geometry} 
\setlength{\parindent}{0pt} % \parindent es la sangria que queremos a principio de párrafo. Puesto a 0pt por defecto.
 % Customización de los headers y footers con fancyhdr https://es.overleaf.com/learn/latex/Headers_and_footers
\usepackage{fancyhdr}
\pagestyle{fancy}
\fancypagestyle{plain}{
  \fancyhf{}
  \fancyfoot[LE,RO]{\thepage}
}
\setlength{\headheight}{15pt} % headheight, separación del encabezado respecto al borde superior de la página
\usepackage{mathtools} % PAquetes para fuentes y símbolos matemáticos
\usepackage{amsmath}
\usepackage{amsfonts}
\usepackage{caption} % Paquete para definir la distancia entre el caption y la tabla
\captionsetup[table]{skip=10pt} % 10pt por defecto
\usepackage{pdfpages} % Paquete para insertar PDFs como la cubierta.
% Configuramos las referencias con biblatex y biber.
% Ejemplo: Bibliografia separada para cada capitulo, estilo de la American Chemical Society, sin DOI, ISBN, URL, etc y con iniciales + apellido para el nombre de los autores
\usepackage[backend=biber,style=chem-acs,refsection=chapter,style=numeric,bibencoding=utf8,doi=false,isbn=false,url=false,eprint=false,date=year,giveninits=true]{biblatex}
\addbibresource{main.bib} % Definimos el archivo .bib donde está la bibliografía
\usepackage{graphicx} % Paquete obligatorio ara insertar imágenes
\usepackage{csquotes}
\usepackage[colorlinks=false,citecolor=black,linkcolor=black,filecolor=black,urlcolor=black]{hyperref} % Referencias en color negro
\graphicspath{ {images/} } % Definimos el directorio matriz donde están las imágenes

% Ir añadiendo los nuevos paquetes que necesitemos aquí

\usepackage{lipsum} % BORRAR Es para generar el texto de prueba Lorem Ipsum

% ---------------------- Fin de la configuración -------------------------------------------

% Comienzo del documento
\begin{document}

% Cubierta
\includepdf[pages={1}]{documents/B5_coberta_CDEgen2017.pdf} 

%Pagina en blanco
\clearpage
\thispagestyle{empty}
\null
\newpage

% Portada
\input{documents/portada}

% Pagina en blanco
\clearpage
\thispagestyle{empty}
\null
\newpage

% Acreditacion
\input{documents/acreditacio}

%Pagina en blanco
\clearpage
\thispagestyle{empty}
\null
\newpage

% Sin header ni footer para los agradecimientos
\fancyfoot{}
\fancyhead{}

\chapter*{Acknowledgements} % Agradecimientos
\pagenumbering{Roman} % Paginas en numeros romanos para el prólogo 
\input{chapters/acknowledgements} % Archivo .tex de los agradecimientos

\tableofcontents % Indice

\fancyfoot[LE,RO]{\thepage} % Pie de pagina, posición del numero de pagina, Left Even (Izda en pares pares) Right Odd (dcha en impares)

% Resumen de la tesis
\chapter*{Overview of the thesis}
\input{chapters/overview}

% Cuerpo de la tesis, cambiar nombre de los capitulos y de los archivos a conveniencia
\chapter{Introduction}
\pagenumbering{arabic} % Paginas en numeros arabigos para el cuerpo de la tesis
\fancyhead[LE]{Chapter \thechapter. Introduction} % Encabezado de pagina, titulo del capitulo a la izda en paginas pares 
\input{chapters/introduction}
\fancyhead{} % Limpia el encabezado para el siguiente capitulo

\chapter{Methodology}
\fancyhead[LE]{Chapter \thechapter. Methodology}
\input{chapters/methodology}
\fancyhead{}

\chapter{Results}
\fancyhead[LE]{Chapter \thechapter. Results}
\input{chapters/results}
\fancyhead{}

\chapter{Conclusions}
\fancyhead[LE]{Chapter \thechapter. Conclusions}
\input{chapters/conclusions}

\thispagestyle{empty}
\null
\newpage
\includepdf[pages={1}]{documents/B5_contra.pdf} 
\end{document}
